\subsection{Connected rough cells in the Wilson seam ensemble}\label{cl:sec:rough-seams}

Write $x_p=N_c^{-1}\ReTr U_p$, so the actual Wilson weight is
$\exp(b\sum_p x_p)$. On a four-dimensional periodic rectangular lattice
with even sides at least four, let $H$ denote its number of unit
four-cells. There are $6H$ plaquettes and $4H$ positively oriented links.
Call a cell rough if at least one of its square faces satisfies
$1-x_p\ge\kappa$. Cells are adjacent if they share a three-dimensional face.

\begin{theorem}[Joint probability of rough cells]\label{cl:thm:rough-cells}
For $\SU(2)$ and $b\ge1$ put $L_b=6+6\log(8b)+4\log(54\pi/25)$.
For $\kappa>0$ and any set $S$ of distinct four-cells,
\[
 {\mathbb P}_b(\text{every cell in }S\text{ is rough})
 \le \min\{1,\exp[(L_b-b\kappa/4)|S|]\}.
\]
In particular $\kappa_b=4[L_b+\log(128b)]/b$ gives an upper bound
$(128b)^{-|S|}$. This threshold tends to zero as $b\to\infty$.
\end{theorem}

\begin{proof}
First disseminate the event from the cells of $S$ to the whole torus by
the chessboard estimate \cite[Theorem~5.8]{Biskup09}. We give the argument
for link variables because adjacent closed cells have common faces.

Cut perpendicular to an axis of length $h$. Let $S_i$ be the pattern of
marked cells in slab $i$. For a tangential pattern $C$, denote by $K_C$
the one-slab integral kernel between its two boundary link fields.
Integrate the normal links with probability Haar measure. Include all
plaquettes with a normal direction, half the action on each boundary,
and the indicators for the cells of $C$. Reflection across the middle
of the slab exchanges the boundary fields, preserves the marked
tangential positions, and inverts the normal links. The real trace of a
plaquette is invariant under this operation up to reversal of its path.
The event uses all faces, so $K_C$ is symmetric. It is bounded and
Hilbert--Schmidt on the compact boundary field space.

The restricted partition integral is the trace of the cyclic product
of the $K_{S_i}$. By the Schatten H\"older inequality it is at most
$\prod_i(\Tr |K_{S_i}|^h)^{1/h}$. Since $h$ is even, each trace here is
$\Tr K_{S_i}^h$, the integral for the pattern $S_i$ repeated in every
slab. This also explains why positivity of the spectrum of a restricted
slab kernel is unnecessary. Divide by the unrestricted partition
function and repeat successively in the other three directions. The
resulting product has one factor per original cell, with exponent $1/H$.
A marked cell contributes the probability that every cell is rough, and
an unmarked cell contributes one. Hence, writing $B_S$ for the joint event,
\[
 {\mathbb P}_b(B_S)\le {\mathbb P}_b(B_{\rm all})^{|S|/H}.
\]

Each plaquette belongs to exactly four four-cells. Thus on $B_{\rm all}$,
$\sum_p(1-x_p)\ge\kappa H/4$, and the restricted partition integral
is at most $e^{6bH-b\kappa H/4}$. For the denominator, identify
$\SU(2)$ with the unit quaternions. Put $\delta=(8b)^{-1/2}$ and restrict
every link to the chord ball of radius $\delta$ about the identity.
Its Haar measure is at least $25\delta^3/(54\pi)$ by the ball estimate
in the proof of Theorem~\ref{cl:thm:frustrated}. Bi-invariance and the
triangle inequality give $\|1-U_p\|\le4\delta$ in quaternion chord length,
hence $1-x_p\le8\delta^2$. Multiplication over the $4H$ links yields
$Z\ge e^{6bH-L_b H}$. Taking the ratio proves the theorem.
\end{proof}

For a fixed compact simple matrix group with a faithful unitary
representation the same proof applies with
$d(g,h)=\|g-h\|_{\rm HS}/\sqrt{N_c}$, so
$1-x_p=d(I,U_p)^2/2$. If $d_G$ is the group dimension, there are constants
$c_G\in(0,1]$, $\delta_G>0$ such that
$\operatorname{Haar}(B_\delta)\ge c_G\delta^{d_G}$ for
$0<\delta\le\delta_G$. Indeed the exponential map has invertible
differential at zero, Haar measure has a positive smooth density in
this chart and the chord length is bounded above by a constant times
the coordinate norm. Integrate over a coordinate ball of the
corresponding radius. Taking $b$ large enough that
$(8b)^{-1/2}\le\delta_G$ gives the theorem with
$L_{b,G}=6+4\log(1/c_G)+2d_G\log(8b)$.
All constants here are for the specified group and representation.

\begin{corollary}[Removing long rough clusters]\label{cl:cor:seam-clusters}
Use the threshold $\kappa_b$ of Theorem~\ref{cl:thm:rough-cells}, or its
group-specific version. Let $T\ge4$ be even and $V=(24T)^3$.
In the weighted matching ensemble of \eqref{cl:eq:seam-disagreement},
form rough clusters on the active Wilson lattice at each corner. Let
$\mathcal G$ be the event that every such cluster has fewer than $T/2$
cells. Then
\[
 \left|p_{\rm diff}-
   {\mathbb P}(\sigma\ne\tau\mid\mathcal G)\right|
 \le {\mathbb P}(\mathcal G^c)
 \le 432 T^4(2b)^{-T/2}.
\]
For $b\ge3$ in the range of the theorem, the last expression is at
most $27648/b^2$ uniformly over all even $T\ge4$.
\end{corollary}

\begin{proof}
A cluster of at least $m$ cells contains a connected set of exactly $m$
containing any given root of the cluster. A rooted spanning tree of such a
set has a depth-first traversal of $2m-2$ steps, with at most eight choices
per step. Each traversal determines its visited set. The joint probability
bound therefore gives
\[
 {\mathbb P}(|\mathcal C(x)|\ge m)
 \le 8^{2m-2}(128b)^{-m}=(2b)^{-m}/64.
\]
Conditional on the two matching signs, the active graph consists of
two Wilson tori of period $T$ or one of period $2T$. In either case it
has $2VT$ potential roots. Sum the last inequality over these roots and
average over the four corners. With $m=T/2$ this gives the stated
coefficient. Along even times the ratio of successive upper bounds is
$(1+2/T)^4/(2b)\le81/(32b)<1$ for $b\ge3$. Substituting $T=4$ gives the
uniform estimate. The event $\mathcal G$ has positive probability since a
neighbourhood of the identity field is included. Finally, for any event
$A$, the total probability formula gives
$|{\mathbb P}(A)-{\mathbb P}(A\mid\mathcal G)|\le{\mathbb P}(\mathcal G^c)$.
\end{proof}

For cuts separated by $T/2$ time steps, $\mathcal G$ excludes a connected
chain of rough cells joining the two cuts. The estimate averages over
the active Wilson fields at the corners. Conditioning on prescribed
coarse links or an exterior field would change the measure in the
chessboard argument. The result gives no such conditional estimate.

The remaining probability in Corollary~\ref{cl:cor:seam-clusters} still
includes long-wavelength gauge fields. For instance, take all links on
a spatial circle of length $n$ in one coordinate direction to be the
constant matrix $\operatorname{diag}(e^{i\theta/n},e^{-i\theta/n})$,
at every transverse site, and all other links equal to the identity.
Every plaquette is exactly flat, while the normalized winding trace is
$\cos\theta$. Local smallness therefore leaves genuine holonomy variables.

There is slack in the initial spectral estimate. In the notation of
Theorem~\ref{cl:thm:seam-disagreement}, $p_{\rm diff}\ge53/134$ gives
$(1+D)^2/(1+E)\le81/53$. Using $E\le D$ gives $D\le28/53<1$.
Using $E\le D^2$ then gives $(7D-2)(2D-7)\ge0$, hence $D\le2/7$.
The same induction yields the same physical gap. Thus, on even cubes
in the theorem's tuning, it suffices that the conditional disagreement
probability exceeds $53/134$ by any positive cutoff-independent margin
as $b\to\infty$. The uniform error in the corollary eventually fits
within that margin. In particular a limiting lower bound of $2/5$
suffices since $2/5-53/134=3/670$. This remaining estimate concerns
the relative weights of the two matchings within $\mathcal G$.

\begin{proposition}[Flat spatial fields with incompatible winding]
\label{cl:prop:seam-winding}
Fix an even $T\ge4$ and the $\SU(2)$ spatial side $n=24T$. For all
sufficiently large $b$ there are open sets of spatial fields of positive
Haar measure on which the conditional probability of opposite seam signs,
given the spatial fields and $\mathcal G$, tends uniformly to zero as
$b\to\infty$. Here $\mathcal G$ is the event of Corollary~\ref{cl:cor:seam-clusters}.
\end{proposition}
\begin{proof}
Take the reference spatial field on sheet zero to be the identity.
On sheet one use $g=\operatorname{diag}(e^{i\pi/n},e^{-i\pi/n})$
in direction one and the identity in both other directions, at every
site and time. Require every spatial link to lie within chord distance
$1/b$ of its reference value. These open sets have positive product Haar
measure for each $b$.

Opposite matching signs require one swapped cut. On each direction-one
row of that cut, let $U_j$ be the active temporal links, with $U_n=U_0$.
For reference endpoints $g_-,g_+$ the $j$th plaquette has chord length
from the identity equal to $d(U_j,g_-U_{j+1}g_+^{-1})$. The map
$u\mapsto g_-u g_+^{-1}$ is an isometry, and its $n$th iterate gives
$-u$, since exactly one of $g_-,g_+$ equals $g$ and $g^n=-I$.
The triangle inequality along the row therefore gives total chord length
at least two. Perturbing the two spatial links in each plaquette costs
at most $2/b$ per plaquette. For $b\ge2n$ the total remains at least one.
Cauchy--Schwarz gives $\sum_{\rm row}(1-x_p)\ge1/(2n)$. The two sheets
have $2n^2$ disjoint rows, giving total deficit at least $n$.

For the straight-straight corner, restrict every active temporal link to
the chord ball of radius $b^{-1/2}$. Each temporal plaquette then has
chord length at most $2b^{-1/2}+2/b$, and each spatial plaquette has
length at most $4/b$. Thus every deficit is at most $8/b<\kappa_b$ for
large $b$, and the restricted corner lies in $\mathcal G$.
There are $2VT$ active temporal links and $12VT$ plaquettes in either
corner geometry. With $h_b$ the Haar mass of the ball used here,
comparison of the two opposite-sign corner weights to this one
straight-straight contribution gives
\[
 {\mathbb P}(\sigma\ne\tau\mid Q_{\rm spatial},\mathcal G)
 \le 2e^{-bn+96VT}h_b^{-2VT}.
\]
Since $h_b\ge(25/(54\pi))b^{-3/2}$, the right side tends to zero at
fixed $n,T$, uniformly on the stated open sets. All inactive link
integrals have mass one.
\end{proof}
This calculation explains why a lower bound uniform in the prescribed
spatial fields cannot finish the argument. The unconditional seam test
requires their weights under the Wilson measure along the physical tuning.

\subsection{A finite matrix for the thermal defect of the interacting theory}
For $\SU(2)$ a complementary estimate controls the tail of the transfer
trace at every coupling. Let $n=3V$ be the number of spatial links.
View their unit quaternions as a vector $Q$ in $\mathbb R^{4n}$, so
$Q\cdot Q'= \sum_e \tfrac12\ReTr(Q_e Q_e'{}^*)$.
Write $\mathbb T=D\Pi C\Pi D$ as in Proposition~\ref{sc:prop:wilson},
where $C(Q,Q')=\exp(b\,Q\cdot Q')$ and $\Pi$ averages vertex gauges.
Define $C_K$ by retaining degrees $0$ through $K$ in this exponential,
and put $\mathbb T_K=D\Pi C_K\Pi D$. Both occurrences of $D$ retain
the complete spatial plaquette action.

\begin{proposition}\label{cl:prop:finite-matrix-seam}
For real $b\ge0$ and an integer $k\ge0$, take $K=\lceil4bn\rceil+k$.
Put $J=\binom{4n+K}{K}$ and
$\varepsilon=4^{-k-1}e^{-(4\log4-5)bn}$. With the exact eigenvalue
ratios indexed from $r_0=1$, for every integer $m\ge1$ one has
\[
 \sum_{j\ge J}r_j^m\le\varepsilon^m.
\]
Writing $\psi_{s,K}(T)=2\log\Tr\mathbb T_K^T-\log\Tr\mathbb T_K^{2T}$,
one also has
\[
 |\psi_s(T)-\psi_{s,K}(T)|\le -2T\log(1-\varepsilon).
\]
\end{proposition}
\begin{proof}
Each kernel $(Q\cdot Q')^j$ is the inner product of the tensor powers of
$Q$ and $Q'$. The Taylor terms are positive operators. Expanding in
monomials shows that the range of $C_K$ has dimension at most $J$.
The positive tail has trace equal to the sum of its diagonal terms:
for each finite-rank monomial kernel this follows by integration, and
the series of these traces converges. Since $Q\cdot Q=n$,
\[
 \Tr(C-C_K)=\sum_{j>K}(bn)^j/j!
 \le 4^{-K-1}e^{4bn}.
\]
Compression by $\Pi$ and multiplication by $D$ preserve operator order.
Cyclicity of the trace and $\|D\|^2\le e^{bn}$ give
$0\le \Tr(\mathbb T-\mathbb T_K)\le 4^{-K-1}e^{5bn}$.
The constant Taylor term in $\mathbb T_K$ is
$|D1\rangle\langle D1|$. Its eigenvalue $\int D^2\,dQ$ is at least
one by Jensen's inequality, since each spatial plaquette has zero Haar
mean. In particular $\lambda_0=\|\mathbb T\|\ge1$.

After division of both operators by the same $\lambda_0$, their ordered
eigenvalues $r_j$ and $u_j$ obey $0\le u_j\le r_j$ by the min--max
principle, and $\sum_j(r_j-u_j)\le\varepsilon$. For $j\ge J$ one has
$u_j=0$, proving the tail estimate. Also
$\|r-u\|_{\ell^m}\le\|r-u\|_{\ell^1}\le\varepsilon$
and $\|r\|_{\ell^m}\ge1$. The reverse triangle inequality therefore gives
$\|u\|_{\ell^m}\ge(1-\varepsilon)\|r\|_{\ell^m}$. The error
in each log trace lies between zero and $-m\log(1-\varepsilon)$.
Taking twice the error at $T$ minus the error at $2T$ proves the
defect estimate.
\end{proof}

The finite matrix can be written explicitly. Index monomials by
multi-indices $\alpha$ of degree at most $K$, and set
$f_\alpha=D\Pi\bigl[(b^{|\alpha|}/\alpha!)^{1/2} Q^\alpha\bigr]$,
with the constant monomial equal to one before applying $D\Pi$.
Then $\mathbb T_K=\sum_\alpha|f_\alpha\rangle\langle f_\alpha|$.
The Gram matrix of the $f_\alpha$ has the same nonzero eigenvalues.
Thus the proposition bounds the contribution lost in this explicit
finite-dimensional approximation. Its dimension increases with the
cutoff, and no bound on the relative weights of the retained eigenvalues
along the tuning follows from the tail estimate.

\subsection{How small a disagreement probability suffices?}
\begin{theorem}[Aspect ratio and the seam threshold]\label{cl:thm:seam-aspect}
Choose integers $R>m\ge2$, a number $0<d<1$, and put
\[
 A_m=\left(\frac{Rm}{R-m}\right)^3,\qquad
 p(d)=\frac{1+d^2}{(1+d)^2+1+d^2}.
\]
Suppose $A_md^m<1$ and $B_m d^{m-1}\le c<1$, where
$B_m=A_m/(1-A_md^m)$. At fixed coupling it suffices to have
$p_{\rm diff}\ge p(d)$ at time $T\ge3$ on a spatial cube of side $RT$.
With $q=c^{1/(m-1)}$, every integer $n\ge RT$ then satisfies
\[
 r_1(n^3)^T\le q,\qquad \Dexc{n^3}{n}\le q^{n/T}.
\]
Under the tuning and continuum-transfer hypotheses of
Theorem~\ref{cl:thm:seam-disagreement}, with $R,m,d,c$ fixed across the
tuning and the test imposed at side $RT$, the physical gap is at least
$-\log(c)/[(m-1)t_0]$.
\end{theorem}
\begin{proof}
Write $D=\sum_{j\ge1}r_j^T$, $E=\sum_{j\ge1}r_j^{2T}$ on the test box.
The probability hypothesis gives
$(1+D)^2/(1+E)\le(1+d)^2/(1+d^2)<2$.
Since $E\le D$, first $D<1$. Since $E\le D^2$ and
$u\mapsto(1+u)^2/(1+u^2)$ is increasing on $[0,1]$, we then have $D\le d$.

Put $T_i=m^iT$ and $D_i=\Dexc{(RT_i)^3}{T_i}$. On the same box the
excited trace at $mT_i$ is at most $D_i^m$. For $RT_i\le n\le RmT_i$,
Theorem~\ref{sc:thm:cube-control} at time $mT_i$ therefore gives
\[
 \Dexc{n^3}{mT_i}\le(1+D_i^m)^{A_m}-1 .
\]
For $0\le x\le d$, the right hand side with $D_i=x$ is at most
$e^{A_mx^m}-1\le A_m x^m/(1-A_m x^m)\le B_m x^m$.
At the endpoint $n=RmT_i$ this proves inductively that
$D_i\le d q^{m^i-1}$, since $B_m d^{m-1}\le q^{m-1}$.
For an intermediate $n$ the very same bound gives
$\Dexc{n^3}{mT_i}\le d q^{m^{i+1}-1}\le q^{m^{i+1}}$,
where $d\le q$ follows from $B_m>1$. Taking the largest excited
eigenvalue gives $r_1(n^3)^T\le q$. The remaining $n-mT_i\ge0$
powers give the second inequality. On the tuning $T\le Nt_0$, the
transfer argument of Theorem~\ref{cl:thm:seam-disagreement} applies
once $L_p\ge Rt_0$.
\end{proof}

On the original cube $R=24$, take $m=18$ and $d=4/9$. Here
$A_m=72^3$, $A_md^{18}<1/5$, $(5/4)A_md^{17}<1/2$, and $p(d)=97/266$.
This gives a sufficient Wilson probability of $97/266$ on the very
same box, with physical gap at least $\log(2)/(17t_0)$ under the
stated hypotheses. On even boxes a fixed positive margin above
$97/266$ in the ensemble conditioned on $\mathcal G$ suffices by
Corollary~\ref{cl:cor:seam-clusters}.

More generally, prescribe any $p_*$ with $1/3<p_*<1/2$. The identity
for the signs and $E\le D$ give $D\le d_*=1/p_*-2<1$ whenever
$p_{\rm diff}\ge p_*$. Choose $m$ large enough that
$16m^3d_*^{m-1}\le1/2$ and take $R=2m$. Then $A_m=8m^3$,
$A_md_*^m<1/2$ and $B_m\le2A_m$. The
induction just proved starts from $D\le d_*$ with $c=1/2$. Thus every
fixed margin above one third suffices on a cube with the corresponding
fixed aspect ratio. In this geometry the same root count as in
Corollary~\ref{cl:cor:seam-clusters} gives
$R^3T^4(2b)^{-T/2}/32$ for the discarded probability on even boxes.
The endpoint illustrates a spectral limitation of this argument: an
abstract transfer system with exactly two equal eigenvalues and
$Z(s,T)=2$ on every box has $p_{\rm diff}=1/3$ and $r_1=1$.
For the interacting Wilson law a uniform lower bound at the required
aspect ratio still has to be proved.
